\Needspace{17\baselineskip}
\SourceLeaf{134}
\noindent Im Übrigen bemerken wir noch, dass, wenn man die Zahlen der ersten Spalte als die doppelten Logarithmen von Tangenten auffasst, die Zahlen der zweiten und dritten Spalte die doppelten Komplemente der Logarithmen der zugehörigen Sinus und Kosinus sind. Mit den gewöhnlichen trigonometrischen Tafeln kann man den beabsichtigten Zweck daher auf genau dieselbe Weise erreichen wie mit den von \textsc{Leonelli} vorgeschlagenen. Die letztere Tafel erspart allerdings die Division und Multiplikation mit 2, die man mit mäßiger Rechenübung leicht im Kopf ausführt. Außerdem liefert sie bei gleicher Zahl von Dezimalstellen die doppelte Genauigkeit. Es ist sonderbar, dass \textsc{Delambre} diese auf der Beobachtung beruhende Art, die trigonometrischen Tafeln zum selben Zweck zu verwenden, nicht erwähnt, obwohl er eine andere, weit weniger bequeme anführt. \textsc{Leonelli} hat ganz recht, sich darüber zu beschweren, dass man ein solches Verfahren mit dem seinen auf eine Stufe stellen wollte.\par\medskip\begin{center}\rule{.24\linewidth}{.3pt}\end{center}

